\chapter{Citations and References}
\label{chap:literature}
Third party ideas must always be identified. Above all, they must be verifiable 
and findable. The technique of citing and referencing serves this purpose.

Different disciplines and degree programs follow specific citation conventions. 
Important information on common citation systems and styles can be found in the 
Writing Center's e-learning courses%
\footnote{\href{https://ilu.th-koeln.de/goto.php?target=cat_52109&client_id=thkilu}{https://ilu.th-koeln.de/goto.php?target=cat\_52109\&client\_id=thkilu}}.

\section{Citations}
Quotation marks at the bottom and top (\enquote{\ldots}) are used to indicate 
text passages that have been copied verbatim.

If your quote already contains a quotation, you must replace the \enquote{double 
quatation marks} in the text with \enquote*{single quotation marks}. See 
also \cref{sec:specialCases}.

\subsection{References}
A source reference must be created for all citations. The source reference is a 
short notation in the citation style which refers to the complete reference in 
the bibliography. References can be made \emph{either} in the continuous text 
(Anglo-American citation style or Harvard style) \emph{or} via a footnote at the 
bottom of the page \emph{or} in an endnote at the end of the entire text. 
Different professional conventions apply here, which must be observed.

If you opt for short references in the text flow or for endnotes, the footnote 
area can be used for comments and for references to passages in your own text.

\emph{Note the positions of the superscript and punctuation marks:}

If the literal quotation itself ends with a period, this is placed before the 
closing quotation mark. The superscript then follows directly afterwards without 
a space. A period is then omitted for the sentence itself. 

If the literal quotation does not end with a period, there are two cases: If the 
quotation is in the middle of your own sentence, the superscript follows 
directly after the quotation marks. If the quotation is at the end of your own 
sentence, first add the quotation marks, then your own period and only then the 
superscript.

To ensure that the source references remain unambiguous even if several short 
titles or year numbers of an author have the same name, the year numbers are 
given an additional small letter.

\subsection{Ibid}
Today, word processing programs make referencing via \emph{ibid} superfluous, as 
you are no longer forced to type out the same information repeatedly. . If you 
nevertheless decide to use this method (e.\,g. on the advice of your examiner), 
we strongly recommend that you only insert ibid. in the final editorial phase, 
as the references are only clear then.

\subsection{Secondary Quotations}
Secondary quotations, are quotations from an author without having checked the 
original source for the context. Secondary citations are only permitted if the 
original source cannot be obtained. In the case of generally accessible 
scientific literature, secondary quotations should be avoided at all costs. If 
it is not possible to compare a quotation with the original text, then you 
should note cited by (followed by the source from which the quotation was 
taken) or cited in. This occurs, for example, when quoting facts from textbooks on the history of science.

\subsection{Citing Figures and Tables}
\label{sec:refFigures}
The same guidelines apply here as for text quotations, i.\,e. there are also 
verbatim and modified reproductions. In the case of a verbatim reproduction, you 
can use an illustration e.\,g. in PowerPoint or a drawing\slash vector program. 
Copyright must be taken into account when scanning; this is only permitted if 
the author or publisher has granted permission~--~usually for a fee. In this 
case, the source reference must be made in the same way as for a literal 
quotation. For your own changes, an addition must be added to the source 
reference (e.\,g. \emph{with minor changes}, \emph{with own calculations}, etc.).

Because graphics are often copied, it is advisable to mark your own graphics or 
illustrations as such

In table and figure captions, the source is always given directly below the 
figure and not in a footnote or endnote. However, the source information does 
not belong in the list of figures and tables. See \cref{sec:captions}.

\section{Bibliography}
The structure of the mandatory bibliography is explained below.

\subsection{Content and Order of the Bibliography}
\label{sec:bib-content}
In the bibliography, all publications used are listed in alphabetical order by 
author's name. The bibliography must include all sources cited in the text, but 
no additional sources. If works are published not only in printed form but also 
electronically, the bibliography should follow the printed version, unless the 
digital document has a fixed document number (DOI). If a letter abbreviation is 
used in the references in the text section instead of the title, this should 
also be indicated in the bibliography.

Several works published by one author within one year must be differentiated. 
For the short titles in the text, the years must therefore be distinguished by 
appending lower case letters. However, the year in the source citation remains 
unchanged without appended letters. The order of a, b, c etc. depends on the 
order of the source references. Whether first names are abbreviated or written 
out in full depends on the conventions of your subject. However, please remain 
consistent.

\subsection{References in the Bibliography}
The order in the notation of the bibliography depends on the document type and 
subject-specific conventions. It is therefore different for monographs than for 
journal articles, and again different in chemistry than in the humanities or 
mathematics. There are also no uniform standards for punctuation.

The data for a bibliography can be found on the title page of a book. This is 
not the book cover, but a printed sheet at the beginning of the book with the 
most important information identifying the work and book. First, the author or 
editor is named together with the title, followed by publication and printing 
information such as the publisher, place and year of publication. Although there 
is a DIN rule for title lettering, it is interpreted very differently.

In the e-learning courses of the Writing Center%
\footnote{\href{https://ilu.th-koeln.de/goto.php?target=cat_52109&client_id=thkilu}{https://ilu.th-koeln.de/goto.php?target=cat\_52109\&client\_id=thkilu}}
you will find further information on the bibliographic information for different 
types of publications. The following bibliography is just one example of a 
possible format.
